\section{Technical lemmas for active-boundary inference}

This supplement records the technical structure underlying the principal
results in the main paper.  Throughout, $p$ denotes observation dimension and
$d$ denotes total parameter dimension.  Vectors and matrices are typeset in
bold; scalar component scores and pairwise contrasts are not.

\subsection{Proof map for the main results}

Theorem~1 is proved in the main text using Lemmas~1.1--1.2 below.
Theorems~2--3 are proved there by the stated uniform law and argmax argument.
Theorems~4--6 use Lemma~1.3 and the displayed bracketing consequence below,
with their final expansions proved in the main text. Theorem~7 has its
conditional-bootstrap proof in the main text. Theorems~8--9 use the boundary-
estimator expansion and Lemma~1.4 below; the bias, variance, Lindeberg, and
plug-in steps are recorded both there and here. Theorem~10 follows from
Theorem~9 and continuous mapping, as proved in the main text. Thus every main
theorem has a same-numbered proof in the manuscript and an explicit
cross-reference here to its supplementary supporting argument.

\subsection{Differentiating the population maximum}

For a fixed active pair $(j,l)$, let
\[
 g(\mathbf{x};\boldsymbol{\theta})
 =q_j(\mathbf{x};\boldsymbol{\theta})
  -q_l(\mathbf{x};\boldsymbol{\theta}),
 \qquad
 \mathbf{s}(\mathbf{x})
 =\nabla_{\boldsymbol{\theta}}
   g(\mathbf{x};\boldsymbol{\theta}^*).
\]

\begin{lemma}[Local positive-part differentiation]
\label{lem:s-positive-part}
Let $\mathcal O$ be an open set on which only components $j$ and $l$ can
attain the maximum.  Under the regular-face and envelope assumptions of the
main paper, for arbitrary $\mathbf{a},\mathbf{b}\in\mathbb R^d$,
\begin{multline}
 \left.
 \frac{\partial^2}{\partial s\,\partial t}
 P_0\!\left[
 \{g(\mathbf{X};\boldsymbol{\theta}^*
       +s\mathbf{a}+t\mathbf{b})\}_+
 \ind\{\mathbf{X}\in\mathcal O\}
 \right]
 \right|_{s=t=0}\\
 =P_0\!\left[
 \ind\{g(\mathbf{X};\boldsymbol{\theta}^*)>0,
          \mathbf{X}\in\mathcal O\}
 \mathbf{a}^\mathsf T
 \nabla_{\boldsymbol{\theta}}^2g
 (\mathbf{X};\boldsymbol{\theta}^*)\mathbf{b}
 \right]\\
 +\int_{\mathcal O\cap\{g=0\}}
 \frac{p_0(\mathbf{x})
 (\mathbf{a}^\mathsf T\mathbf{s}(\mathbf{x}))
 (\mathbf{b}^\mathsf T\mathbf{s}(\mathbf{x}))}
 {\|\nabla_{\mathbf{x}}g
       (\mathbf{x};\boldsymbol{\theta}^*)\|}
 \,d\mathcal H^{p-1}(\mathbf{x}).
 \label{eq:s-positive-part}
\end{multline}
\end{lemma}

\begin{proof}
Choose a nonnegative, symmetric $C^\infty$ density $K_0$ supported on
$[-1,1]$ and define
\[
 \rho_\varepsilon(v)
 =\int (v-\varepsilon z)_+K_0(z)\,dz.
\]
Then $|\rho_\varepsilon(v)-v_+|\leq\varepsilon$ and
\[
 \rho_\varepsilon'(v)=\int_{-1}^1\ind\{v>\varepsilon z\}K_0(z)\,dz,
 \qquad \rho_\varepsilon''(v)=\varepsilon^{-1}K_0(v/\varepsilon).
\]
Put $g_{s,t}(\mathbf x)=g(\mathbf x;\boldsymbol\theta^*+s\mathbf a+t\mathbf b)$
and $F_\varepsilon(s,t)=P_0[\rho_\varepsilon(g_{s,t}(\mathbf X))
\ind\{\mathbf X\in\mathcal O\}]$. For fixed $\varepsilon$, the chain rule gives
\begin{align*}
 \partial_{st}F_\varepsilon(0,0)
 &=P_0[\rho_\varepsilon'(g)\mathbf a^\mathsf T
 \nabla_{\boldsymbol\theta}^2g\,\mathbf b\ind\{\mathbf X\in\mathcal O\}]\\
 &\quad+P_0[\varepsilon^{-1}K_0(g/\varepsilon)
 (\mathbf a^\mathsf T\mathbf s)(\mathbf b^\mathsf T\mathbf s)
 \ind\{\mathbf X\in\mathcal O\}].
\end{align*}
The coarea formula writes the second term as
\[
 \int_{-1}^1K_0(z)\int_{\mathcal O\cap\{g=\varepsilon z\}}
 \frac{p_0(\mathbf x)(\mathbf a^\mathsf T\mathbf s(\mathbf x))
 (\mathbf b^\mathsf T\mathbf s(\mathbf x))}
 {\|\nabla_{\mathbf x}g(\mathbf x;\boldsymbol\theta^*)\|}
 \,d\mathcal H^{p-1}(\mathbf x)\,dz.
\]
Surface-moment continuity and domination give the boundary integral as
$\varepsilon\downarrow0$. Since $0\leq\rho_\varepsilon'\leq1$ and it converges
to $\ind\{v>0\}$ off zero, dominated convergence gives the interior term.
The uniform approximation bound and derivative envelopes justify commuting
smoothing with both directional difference quotients, so the limit is the
mixed derivative of the unsmoothed positive part \citep{evans2015measure}.
\end{proof}

\begin{lemma}[Assembly over active strata]
\label{lem:s-strata}
Suppose the active faces are locally finite and multiple-tie junctions have
$\mathcal H^{p-1}$-measure zero.  Then a smooth partition of unity can be
chosen so that every partition element intersects either a unique-winner
region or the relative interior of exactly one active face.  Summing
Lemma~\ref{lem:s-positive-part} over this partition produces
\[
 \nabla_{\boldsymbol{\theta}}^2M(\boldsymbol{\theta}^*)
 =-\mathbf{I}_c+\mathbf{B}.
\]
Inactive equality sets contribute zero.
\end{lemma}

\begin{proof}
Remove a decreasing open neighbourhood of the multiple-tie junctions.  On the
remaining closed set, regularity and local finiteness give an open cover by
unique-winner and single-face charts.  A subordinate smooth partition of
unity exists.  Lemma~\ref{lem:s-positive-part} applies on each face chart;
ordinary dominated differentiation applies on unique-winner charts.  Let the
junction neighbourhood shrink.  Its surface contribution vanishes by the
Hausdorff-measure assumption and the surface envelope. At an inactive point of
$q_j=q_l$, a third score $q_m$ strictly dominates both. Continuity preserves
that inequality on a neighbourhood and for nearby parameters. Crossing
$g_{jl}=0$ therefore cannot change the winner, so the local maximum has no
corresponding positive-part kink or delta mass. A locally finite cover applies
this argument to the entire inactive stratum.
\end{proof}

\subsection{Localized empirical-process control}

The margin condition is used precisely to control observations whose winner
can change under a displacement $\mathbf u$. Local Lipschitz continuity makes
a switch possible only when an active contrast lies within
$O(\|\mathbf u\|)$ of zero, and the margin bound makes this tube's probability
$O(\|\mathbf u\|)$. That factor yields the $\|\mathbf u\|^{3/2}$
$L_2(P_0)$ remainder below; it is not needed on smooth unique-winner regions.

Let
\[
 r_{\mathbf{u}}(\mathbf{x})
 =m_{\boldsymbol{\theta}^*+\mathbf{u}}(\mathbf{x})
  -m_{\boldsymbol{\theta}^*}(\mathbf{x})
  -\mathbf{u}^\mathsf T\boldsymbol{\psi}^*(\mathbf{x}),
 \qquad \|\mathbf{u}\|\leq\delta.
\]

\begin{lemma}[Envelope and entropy of the local remainder class]
\label{lem:s-entropy}
Under the local Lipschitz-envelope and boundary-margin conditions, there are
constants $C_1,C_2$ such that
\begin{align}
 \sup_{\|\mathbf{u}\|\leq\delta}
 \|r_{\mathbf{u}}\|_{P_0,2}
 &\leq C_1(\delta^{3/2}+\delta^2),                         \label{eq:s-l2}\\
 \log N_{[]}
 \left(\varepsilon\|F_\delta\|_{P_0,2},
       \{r_{\mathbf{u}}:\|\mathbf{u}\|\leq\delta\},
       L_2(P_0)\right)
 &\leq C_2d\log(C_2/\varepsilon),                         \label{eq:s-entropy}
\end{align}
for $0<\varepsilon<1$, where
$\|F_\delta\|_{P_0,2}=O(\delta^{3/2}+\delta^2)$.
\end{lemma}

\begin{proof}
Decompose the sample space into the event on which the winner is unchanged and
the union of active tubes of width $O(\delta L)$.  Taylor's theorem bounds the
first contribution by $O(\delta^2L)$; the second is bounded by
$O(\delta L)$ and has probability $O(\delta)$.  This proves
\eqref{eq:s-l2}, after truncation and use of $P_0L^{2+\eta}<\infty$.

Cover the Euclidean $d$-ball of radius $\delta$ by
$O(\varepsilon^{-d})$ balls.  The Lipschitz score envelope and the same tube
decomposition create upper and lower brackets around the function indexed by
each center.  Their $L_2(P_0)$ widths are bounded by a constant multiple of
$\varepsilon\|F_\delta\|_{P_0,2}$, which gives
\eqref{eq:s-entropy}.
\end{proof}

The bracketing integral implied by \eqref{eq:s-entropy} is finite.  The
bracketing maximal inequality therefore gives, for fixed $R$,
\[
 \sup_{\|\mathbf{u}\|\leq R}
 \left|\sqrt n\,\mathbb G_n
 r_{\mathbf{u}/\sqrt n}\right|=o_p(1),
\]
which is the nonlinear remainder required by the local empirical-process
expansion in the main paper \citep{van1996weak}.

\subsection{Boundary-estimator expansion}

For each active pair, define the matrix-valued surface moment
$\boldsymbol{\Gamma}_{jl}(v)$ as in the main paper.  By the coarea formula,
the oracle expectation is exactly
\begin{equation}
 E\widetilde{\mathbf{B}}_n
 =\sum_{j<l}\int K(t)\boldsymbol{\Gamma}_{jl}(ht)\,dt.
 \label{eq:s-oracle-expectation}
\end{equation}
A componentwise Taylor expansion, symmetry of $K$, and
$\int K=1$ give the stated $O(h^2)$ bias.  Applying the same identity to the
outer product of vectorized summands gives
\begin{equation}
 nh\,\operatorname{Var}
 \{\operatorname{vec}(\widetilde{\mathbf{B}}_n)\}
 \longrightarrow\boldsymbol{\Omega}_{\mathbf{B}}.
 \label{eq:s-oracle-variance}
\end{equation}
Compact support and the fourth-order surface envelope imply the Lindeberg
condition for the triangular array.

\begin{lemma}[Uniform parameter substitution]
\label{lem:s-plugin}
Let $\overline{\boldsymbol{\theta}}_n$ satisfy
$\|\overline{\boldsymbol{\theta}}_n-\boldsymbol{\theta}^*\|=o_p(h)$.
Under the kernel, derivative-envelope, and junction-neighbourhood assumptions,
uniformly over such sequences,
\begin{equation}
 \|\widehat{\mathbf{B}}_n(\overline{\boldsymbol{\theta}}_n)
   -\widehat{\mathbf{B}}_n(\boldsymbol{\theta}^*)\|_F
 =O_p\!\left(
 \frac{\|\overline{\boldsymbol{\theta}}_n
             -\boldsymbol{\theta}^*\|}{h}+\zeta_h
 \right).
 \label{eq:s-plugin}
\end{equation}
\end{lemma}

\begin{proof}
Split the difference into changes in the kernel argument, derivative outer
product, and active-pair indicator.  Lipschitz continuity of $K$ introduces
the factor $h^{-1}$ in the first term.  Uniform derivative envelopes control
the second.  The active indicator can change only if a third component is
within $O(h+\|\overline{\boldsymbol{\theta}}_n-
\boldsymbol{\theta}^*\|)$ of the tied pair; the normalized contribution of
these junction neighbourhoods is bounded by $\zeta_h$ plus the first term.
The empirical deviations obey the same bounds by the localized maximal
inequality.
\end{proof}

Combining \eqref{eq:s-oracle-expectation},
\eqref{eq:s-oracle-variance}, and Lemma~\ref{lem:s-plugin} proves both
consistency and the plug-in rate stated in the main paper.
